\documentclass[11pt,a4paper]{article}
\usepackage[margin=23mm]{geometry}
\usepackage[T1]{fontenc}
\usepackage{lmodern,amsmath,amssymb,amsthm,microtype}
\usepackage[hidelinks]{hyperref}
\setlength{\parindent}{0pt}
\setlength{\parskip}{0.45em}
\setlength{\emergencystretch}{1em}
\DeclareMathOperator{\Supp}{Supp}
\newtheorem*{lemma}{Lemma}
\hypersetup{pdftitle={Very exceptional negativity}}
\pagestyle{empty}
\begin{document}
\begin{center}\large\bfseries Very exceptional negativity\end{center}
We adapt the critical-coefficient argument in the
\href{https://username1380068.github.io/Note/notes/negativity.html}{negativity note}
to \href{https://www.numdam.org/article/PMIHES_2012__115__325_0.pdf}{Birkar's Lemma~3.3}.

\begin{lemma}
Let $f:X\to Z$ be a projective contraction of normal varieties over an
uncountable algebraically closed field $k$,
and let $D=D^+-D^-$ be an $\mathbb R$-Cartier divisor, where $D^\pm\ge0$
have no common components. Suppose that $D^-$ is very exceptional over $Z$
and, for every component $S$ of $D^-$, $-D|_S$ is nef on the very general
curves of $S/Z$. Then $D\ge0$.
\end{lemma}

\begin{proof}
Suppose $D^-\ne0$. Shrinking $Z$, we may assume that $Z$ is affine and all
components $D_i$ of $D^-$ map onto the same irreducible closed subset
$P\subsetneq Z$. Choose an effective Cartier divisor $A$ containing $P$, and set
\[
 e=\min\{t\ge0:D+t f^*A\ge0\},\qquad
 G=D+e f^*A,\qquad
 S=\bigcup_{\operatorname{coeff}_{D_i}G=0}D_i.
\]
Then $e>0$, $S\ne\varnothing$, no component of $S$ is contained in
$\Supp G$, and $\Supp f^*A\subset S\cup\Supp G$.

Choose a prime component $R$ of $A$ containing $P$. Very exceptionality
gives a prime divisor $Q$ with $f(Q)=R$ and $Q\not\subset\Supp D^-$;
hence $Q\subset\Supp G$. Every fibre over $P$ therefore meets both $S$
and $\Supp G$, and is covered by their union. Connectedness gives
\[
 P=\bigcup_{D_i\subset S}f(D_i\cap\Supp G).
\]
The images are closed, so irreducibility of $P$ yields a component
$D_i\subset S$ such that $F:=D_i\cap\Supp G$ dominates $P$.

Put $d=\dim D_i-\dim P$. We have
$d=(\dim X-\dim Z)+\operatorname{codim}_ZP-1\ge1$:
if $\dim X=\dim Z$, then $f$ is birational and very exceptionality gives
$\operatorname{codim}_ZP\ge2$.

Here the very-general-nef hypothesis means that
\[
\begin{gathered}
 \exists(\Delta_j)_{j\ge1},\quad \Delta_j\subsetneq D_i\text{ closed},
 \quad\forall\text{ integral curves }C\subset D_i:\\
 \dim f(C)=0,\quad C\not\subset\bigcup_j\Delta_j
 \quad\Longrightarrow\quad D\cdot C\le0.
\end{gathered}
\tag{VG}
\]
To contradict this, we show that for \emph{every} countable family of proper
closed subsets $(\Delta_j)_{j\ge1}$ of $D_i$, there is an integral curve
$C\subset D_i$ such that
\[
 \dim f(C)=0,\qquad C\not\subset\bigcup_j\Delta_j,
 \qquad D\cdot C>0.
\tag{$*$}
\]
Fix any such family. Write $W_z=W\cap f^{-1}(z)$ for $W\subset D_i$.
For a very general $z\in P$, the fibre
$(D_i)_z$ is pure of dimension $d$, while $F_z$ and every $(\Delta_j)_z$
have dimension at most $d-1$. Since $F_z\ne\varnothing$, some irreducible
component $T$ of $(D_i)_z$ meets $F$ but is contained in neither $\Supp G$
nor any $\Delta_j$. A very general ample complete-intersection curve
$C\subset T$ (or $C=T$ if $d=1$) is contained in none of these bad sets
and satisfies $G\cdot C>0$, since $G|_T$ is nonzero and effective.
As $f(C)$ is a point, we obtain
\[
 D\cdot C=(G-e f^*A)\cdot C=G\cdot C>0.
\]
This proves $(*)$, contradicting (VG).
\end{proof}
\end{document}
